\section{Introduction}
\label{sec:intro}
Multi-task benchmarks are the standard way to compare machine-learning models: a leaderboard is a table of scores of many models on many sub-datasets, summarised by an aggregation rule, most often the (weighted) mean. Ranks influence adoption, publication and funding, so developers have a reason to act strategically. One such action is \emph{targeted investment}: training the favoured model on data related to some of the sub-datasets so that its scores there rise. Public discussion of leaderboards and contamination~\citep{singh2025illusion,zhou2023cheater,yang2023rephrased} documents the incentive, but not how much effort is needed to reach the top, nor how that depends on the way scores are aggregated.

We study exactly that. \emph{How much targeted investment does it take for a given model to become first, and how does this vary with the aggregation rule?} We take the question to be \emph{conditional by design}: the answer depends on a modelling assumption (how far a score can be raised, our cap scenarios), on the rule, and on the target's starting position. We therefore report every result under a named cap, with three cost measures that are never merged, and with the rate at which first place is attainable at all.

\paragraph{Contributions.}
(1) A formal minimum-investment problem on a score matrix, with explicit cap scenarios, cost measures, margins and a set of rules (Section~\ref{sec:method}); the link to bribery in computational social choice is made explicit in Section~\ref{sec:related}.
(2) Exact solvers for score rules and, by mixed-integer programming, for Borda, plurality and the Condorcet-winner requirement, tested against brute force, together with closed forms for the mean, the minimum and the odd-$n$ median that we verify on every target of two public leaderboards.
(3) An empirical study on HELM MMLU v1.13.0 (79 models $\times$ 57 subjects) and the Open LLM Leaderboard (4496 models $\times$ 6 tasks), with bootstrap intervals over targets and robustness analyses.

\paragraph{Findings (each conditional on the stated cap and data).}
(i)~On HELM, for the ten challengers ranked 2--11 under the mean, reaching first place costs the mean rule a median of 162 points in total over 4 subjects under cap \emph{ceil}; the median rule, the minimum rule and plurality need only 3--6\% of that in total points (Section~\ref{sec:exp}). (ii)~This ordering of rules by cost hardly changes across six cap scenarios, but \emph{feasibility} does: tight caps make first place unattainable for some targets under the minimum rule and the power means. (iii)~The breadth that order-statistic rules force in theory (a majority of tasks for a target below the rival median everywhere) is only partly realised for targets near the top: with the same target for every rule, $K$ under the median rule grows with the number of tasks $n$ (log--log slope $0.54$, 0.51--0.56; median $K=6$ of 48 subjects), but far below $\lfloor n/2\rfloor+1$ and below the mean ($0.84$; $K=13$ once per-task headroom is capped), and $K$ under the minimum is nearly flat ($0.16$). (iv)~The cheap rules have the largest relative spread of cost under evaluation noise. (v)~On the six-task Open LLM Leaderboard the score scale matters for score rules in paired comparisons over targets: the median rule's ratio to the mean rule rises from $0.13$ to $0.84$ (ranks 11--30) and from $0.34$ to $0.78$ (ranks 31--100) under raw instead of normalised scores, but not beyond rank 100, and the minimum is dearer under raw scores in every rank group; rank rules are invariant. We also report what we did not complete (Section~\ref{sec:limits}).
